Ti--Al based materials have attracted significant attention in engineering due to their remarkable mechanical and chemical properties. The intermetallic alloy titanium aluminide (TiAl), for instance, is highly sought after in the aviation and automotive industries for its low density, high strength, and immunity to corrosion and oxidation.\cite{emiraliouglu2022additive} In addition, composite materials such as Ti/TiAl\textsubscript{3} and Al/TiAl\textsubscript{3} compound metal-metallic composites, which are based on Ti and Al, have also gained interest in the aerospace and defense industries for their high yield strength at elevated temperatures, high energy absorption capability, dimensional stability, and corrosion/oxidation resistance.\cite{zhang2022formation,ma2020microstructure} 

Nevertheless, fabrication of Ti--Al based materials can sometimes be challenging. Ti--Al intermetallics such as TiAl and TiAl\textsubscript{3} are brittle at room temperature, making them difficult to produce by conventional methods such as machining or shaping.\cite{emiraliouglu2022additive,dilip2017novel} The production of Ti--Al based materials is often achieved through various methods such as powder metallurgy for alloys,\cite{emiraliouglu2022additive} rolling followed by heat treatment or hot pressing for laminated structures,\cite{zhang2022formation,ma2020microstructure} and additive manufacturing.\cite{emiraliouglu2022additive} One of the key characteristics that these methods have in common is the use of elevated temperatures at specific stages. Notably, higher temperatures allow for increased diffusion of atoms.\cite{thiyaneshwaran2018nucleation} The diffusion at the interface between Ti and Al play a crucial role in the formation of various microstructures, including the Kirkendall voids.\cite{thiyaneshwaran2018nucleation} Such voids are the result of the Kirkendall effect, which is the motion of the interface between two metals that occurs as a consequence of the difference in diffusion rates of the metal atoms.\cite{wang2013hollow,el2015kirkendall} Presence of these voids is not desired in cases where the mechanical strength and stability of the metal part is crucial.\cite{sun2020fracture}

There has been many experimental efforts to characterize the diffusion dynamics of Ti--Al binary systems. It is well established that TiAl\textsubscript{3} is the primary intermetallic phase that forms in the Ti/Al diffusion couple.\cite{mirjalili2013kinetics,zhang2022formation} A more recent study has revealed that in a Ti/Al binary diffusion couple, the diffusion of Al towards Ti is dominant, resulting in the concentrated nucleation of TiAl\textsubscript{3} at the Ti/TiAl\textsubscript{3} interface rather than the Al/TiAl\textsubscript{3} interface.\cite{thiyaneshwaran2018nucleation} However, current studies are primarily based on empirical relations such as the widely used model in which the thickness of reaction layer $x$ is linear in diffusion time $t$ raised to a kinetic exponent $n$, i.e. $x = Kt^n$, where $K$ is the rate constant.\cite{xu2006growth,mirjalili2013kinetics,zhang2022formation} 
These models are powerful tools when explaining and predicting regular systems such as laminated composites, but may not be easy to use when it comes to systems with complex interfaces and microstructures, such as those found in materials produced through additive manufacturing. 
% In such systems, detailed atomistic-level insights is more direct and may provide more insights. 
\cref{fig:kirkendall_experiment} illustrates the cross-sectional views of one of our previous attempts to create Al/TiAl\textsubscript{3} composite using direct energy deposition (DED). The complex nature of the system is compounded by the presence of dendrites, dispersed pure Ti particles, and the non-uniform concentration gradient of Ti and Al. The images were captured using scanning electron microscope (SEM) and show the composite after 0 hours, 6 hours, and 12 hours of heat treatment. As the heat treatment time increases, the Kirkendall effect becomes more significant, as observed by the formation of voids around the Ti particles and intermetallics at the Ti/Al interface. These microstructural changes are a result of the difference in atomic diffusion rates between the titanium and aluminum components of the alloy, and are a clear demonstration of the Kirkendall effect in action. This highlights the importance of understanding the diffusion dynamics in the design and manufactruing of new alloys and composites. 

\begin{figure}[htbp]
    \centering
    \includegraphics[width=0.75\textwidth]{figs/kirkendall_effect_in_AM/kirkendall_experiment.png}
    \caption{Cross-sectional views of a Ti--Al metal composite manufactured by additive manufacturing. The light regions in the image correspond to pure titanium, the gray regions correspond to TiAl\textsubscript{3} in the shape of dendrites that formed during manufacturing and layers that formed during heat treatment, and the dark regions correspond to pure aluminum. The images were taken using scanning electron microscope (SEM) and show the composite after 0 hours of heat treatment (top), 6 hours of heat treatment (middle), and 12 hours of heat treatment (bottom). To the right of each row, a zoom-in of a representative region of the corresponding view is shown to highlight the details of the microstructures near the Ti--Al interface.}
    \label{fig:kirkendall_experiment}
\end{figure}

Molecular Dynamics (MD) simulations are a proven technique for understanding the atomic-level dynamics of materials, including alloys. Several previous studies have demonstrated the efficacy of MD in simulating the formation of Ti--Al intermetallics.\cite{kiselev2014molecular,kiselev2016modeling,levchenko2012molecular} Moreover, a recent MD study combined with experiments on the reactivity of the Ti--Al system have identified an exothermic atom substitution process at the Ti/Al interface, where liquid Al atoms replace solid Ti atoms, leading to continued dissolution from areas where multiple Ti atoms have been substituted in close proximity.\cite{bizot2020reactivity} Another recent study uses MD to investigate the self-diffusion in liquid and solid alloys of the Ti--Al system. The results show that the activation energy of self-diffusion depends on the concentration of components, and is smaller in disordered alloys than in ordered alloys.\cite{poletaev2021self} MD simulations are a powerful tool for understanding physical processes at the atomic level, including the vacancy diffusion mechanism.\cite{poletaev2021self} However, traditional MD analysis often relies on researchers observing and interpreting a few infrequent events in simulation movies. Such analysis can lead to a lack of statistical relevance for the events observed or to potentially miss crucial but rare events.\cite{prinz2011markov}

Markov State Models (MSM) provide a data-driven approach for analyzing molecular dynamics (MD) simulation data. MSM partitions the dynamics of atoms or molecules into discrete states and calculates the transition probabilities based on either theory or data, thus producing a model of the dynamics that is free from the subjectivity of the analyst.\cite{prinz2011markov} Also, it is claimed that all valid equilibrium MD can be explained through the concept of relaxation processes that are represented by the eigenfunctions of the dynamical operator, making MSM a generalizable approach for MD data analysis.\cite{prinz2011markov} The development of the Variational Approach for Markov Processes (VAMP) offers a mathematical framework for optimizing Markov State Models (MSM) with time series data.\cite{wu2020variational} In a follow-up work, the theory of VAMP is implemented in deep neural networks to learn the transition rate between molecular structures.\cite{mardt2018vampnets} Furthermore, the recently developed Graph Dynamical Networks (GDyNets) uses a graph convolutional neural network architecture that incorporates the graph representation of particle systems and the VAMP theory in learning atomic-scale dynamics in general phases and environments from MD simulations.\cite{xie2019graph}

In our study, we first employ MD to scrutinize the behavior of Ti and Al atoms within TiAl\textsubscript{3} grain boundaries at elevated temperatures, thereby replicating the conditions of experimental heat treatment. Through our computational data, we explore the kinetics involved in the formation of TiAl\textsubscript{3} during this heat treatment process. Subsequently, we use MSM to delve deeper into the dynamics of Al atoms in the Ti/Al amalgamation, taking into account the presence of Ti and Al substrates. By characterizing distinct dynamic states, mapping their spatial distributions, and determining the varied timescales associated with state transitions, we are able to provide a comprehensive analysis of the interfacial dynamics of Ti--Al systems, and therefore, shed light on the role of the dynamics in the formation of TiAl\textsubscript{3} during the manufacturing process.